% Introduction. Plain text; empty lines start a new paragraph.

% Optional: replace the default purpose text.
% \Purpose{This document ...}

\Scope{
  In scope: the component \TheComponent\ with its three units in the series
  configuration: its provided interfaces, the interaction of its units, its
  error handling and its resource budgets.

  Out of scope: the internals of each unit (verified by their unit test
  specifications), the NvM, Dem and IoHwAb modules (basic software), and the
  heater power stage (system integration test).
}

% \DocumentReference{key}{attribute = value, ...} -- listed as internal and
% external documents, labelled with id or [R1], [R2], ...; cite with \DocRef{key}
% (one attribute: \DocRef[issue]{key}). Every listed document must be cited.
% Attributes: id (label instead of R<n>), title, number (document number),
% issue, date, scope (internal | external), url, classification (public,
% internal, confidential, strictly confidential). Required: title; for internal
% documents also number and issue. Values with , or = in braces, a url always
% in braces, % in it as \%.
\DocumentReference{srs}{
  title = Software Requirements Specification,
  number = SRS-SPC-001, issue = 1.3, date = 2026-09-14,
  classification = confidential,
  url = {https://plm.example.com/windchill/app/#ptc1/tcomp/infoPage?oid=VR\%3Awt.doc.WTDocument\%3A4711&u8=1},
}
\DocumentReference{arch}{
  title = Software Architecture Description,
  number = SAD-HC-002, issue = 2.0, date = 2026-10-20,
  classification = confidential,
  url = {https://plm.example.com/windchill/app/#ptc1/tcomp/infoPage?oid=VR\%3Awt.doc.WTDocument\%3A4720&u8=1},
}
\DocumentReference{suts-spc}{
  title = Software Unit Test Specification setpoint_ctrl,
  number = SUTS-SPC-001, issue = 0.2,
  classification = internal,
}
\DocumentReference{suts-tf}{
  title = Software Unit Test Specification temp_filter,
  number = SUTS-TF-001, issue = 1.0,
  classification = internal,
}
\DocumentReference{suts-pi}{
  title = Software Unit Test Specification heat_pi,
  number = SUTS-PI-001, issue = 1.1,
  classification = internal,
}
\DocumentReference{plan}{
  id = SVP, title = Software Verification Plan,
  number = SVP-PRJ-002, issue = 2.0,
  classification = internal,
}
\DocumentReference{29119}{
  title = ISO/IEC/IEEE 29119-3 -- Test documentation,
  issue = 2021, scope = external,
  url = {https://www.iso.org/standard/79429.html},
}
\DocumentReference{aspice}{
  title = {Automotive SPICE PAM -- SWE.2, SWE.5},
  issue = 4.0, scope = external,
  classification = public,
}
\DocumentReference{26262}{
  title = ISO 26262-6 -- Clauses 7 and 10,
  issue = 2018, scope = external,
}

% Stakeholders and use cases (optional; usually taken over from the SRS).
% {use case}{actor, actor} -- actors are collected automatically.
\UseCaseNote{System-level use cases the component contributes to, from \DocRef{srs}.}
\UseCase{Set cabin temperature}{Vehicle occupant}
\UseCase{Heat the cabin to the setpoint}{Vehicle occupant}
\UseCase{Read controller state}{Service technician, Diagnostic tester}

% {term}{definition}
\Term{CUT}{Component Under Test}
\Term{TC}{Test Case}
\Term{IS}{Integration Step}
\Term{WCET}{Worst-Case Execution Time}
\Term{NvM, Dem, IoHwAb}{Non-volatile Memory manager, Diagnostic event manager, I/O Hardware Abstraction}
